\section{Discussion}\label{sec:discussion}

\subsection{The picture}

Our current picture of rich-regime gradient descent on these targets \status{Heuristic}, summarizing the evidence of \cref{sec:networks,sec:model}:
\[
\text{trained network}\;\approx\;\underbrace{\text{random-feature kernel background}}_{\text{fits the sample noise; decides which examples are wrong}}\;+\;\underbrace{\text{a few grown features}}_{\text{decide whether and when it generalizes}}.
\]
Features are added greedily. A candidate is accepted once its correlation with the current residual beats the sampling noise of the \emph{continuous} search that would find it. That noise depends on how many inputs the feature needs at once that the circuit does not yet use, or, for a feature built from grown ones, how many existing features it must combine at once (\cref{sec:composition}); it grows polynomially in $d$. It is not the best-of-$M$ noise of a union bound over a finite library, so the cost is not a description length.

\subsection{Relation to NNBayes $\approx$ Occam}

\begin{itemize}
\item \textbf{The same answer to \emph{what} is learned.} Once gradient descent generalizes, the network is functionally close to a small circuit, as the distillation certificates show. So does the Occam picture predict \cite{nngen}.
\item \textbf{Different answers to \emph{how many samples}.}
  \begin{itemize}
  \item Occam pays the description length, about $k\log d$ for a $k$-parity. The $\gamma\ge1$ posterior of \cite{nngen} is proved to need only polynomially many samples; whether it pays only the description length is open.
  \item Gradient descent pays the noise cost of the most expensive move along its greedy path, about $d^{k-1}$ for a move that needs $k$ new inputs at once.
  \end{itemize}
\item \textbf{The gap grows with leap:} the median ratio is $8\times$, $17\times$, $36\times$ and $85\times$ at leaps 1–4, at $d=30$.
\end{itemize}
A conjecture suggested by the predictor, not shown here \status{Conjecture}: the two sample complexities are close for targets with a cheap greedy path, i.e.\ low-leap compositions with large residual correlations at every step. They differ by about $d^{k-1}/(k\log d)$ when a leap of size $k$ is unavoidable. For fixed $k$ that gap is itself polynomial; it becomes super-polynomial only if $k$ grows with $d$.

\subsection{What is in the literature, and what is new}

\paragraph{Known.} The qualitative picture is largely in the literature:
\begin{itemize}
\item the saddle-to-saddle / incremental-learning view of rich-regime gradient descent \cite{pesme2023saddle,jacot2021saddle,li2021greedy,kunin2025agf,bantzis2025saddle};
\item the leap / CSQ account of which sparse targets online SGD learns, and how fast \cite{abbe2022msp,abbe2023leap}, and its analogue for single-index models, the information exponent \cite{benarous2021online};
\item the equivalence of mutation-and-selection with CSQ learning \cite{feldman2008evolvability}.
\end{itemize}

\paragraph{New here, as far as we know.} The contributions are quantitative and empirical:
\begin{enumerate}
\item a calibrated discrete model and a two-constant closed form that predict per-target sample complexity of multi-epoch, fixed-sample training out of sample, including a pre-registered prediction at $d=120$;
\item ablations showing that the finite-sample noise level of the direction search, not the race dynamics, is the operative ingredient;
\item the kernel-background decomposition of the networks' mistakes;
\item the specific theorems of \cref{app:discrete,app:networks}: the weight-free rank bound for quadratic networks, and the early-phase results for ReLU networks;
\item the empirical finding that, for networks trained on a fixed sample, distillation certificates are the only tight bounds we found.
\end{enumerate}

\subsection{Leap counts supports, not structure}\label{sec:leapcritique}

Leap complexity is a function of the target's Fourier \emph{supports} only. Once the inputs of a monomial have been found, any further function of them is free in $d$. Leap therefore cannot see whether a high-order term can be built from intermediate features that are useful in their own right. Our move classification in \eqref{eq:noiselaw} inherits part of this blindness, since $k$ counts new inputs. But the correlation $\rho$ of the best available gate depends on which features already exist, so the noise law can in principle see such structure. An early version of this work proposed a ``circuit leap'' complexity. It orders a circuit's gates so that each insertion is a positive-utility move, and counts the largest number of gates that must be inserted at once. \Cref{sec:composition} tests the question directly.

\subsection{Limitations}
\begin{itemize}
\item \textbf{One setting.} All sample-complexity experiments use one architecture family (a $\mu$P two-hidden-layer ReLU MLP), mostly width 128 and Adam, square loss, uniform Boolean inputs, $d\le120$, and interactions of at most 4 bits. The theory checks use small two-layer ReLU, quadratic and diagonal linear networks.
\item \textbf{Width.} The 4-parity's width dependence is not explained by any model here (\cref{sec:width}).
\item \textbf{The noise law.} Its form for 3 or more simultaneous inputs is a fit, not a derivation. It misses the staircase, multiplexer and MAJ of parities by up to $2.1\times$, prices the reuse move of the composition test several times too cheaply (\cref{sec:composition}), and overpredicts the 4-parity at $d=60$.
\item \textbf{Training budget.} The networks' $\nstar$ is still falling at 20000 steps, by about 0.17 in $\log_2$ per doubling of training. Part of every match is a match at this training budget.
\item \textbf{Verification.} The theorems were checked by AI referees, not by humans (\cref{app:provenance}).
\end{itemize}

\subsection{Open problems}
\begin{enumerate}
\item \textbf{Prove the noise law} for a two-layer network trained on a fixed sample from small initialization: an order-$k$ feature escapes the plateau once $n\gtrsim d^{\max(1,k-1)}/\rho^2$, up to logarithms, with a matching lower bound for $k\le3$. For $k\ge4$ the networks scale below $d^{k-1}$, and the right exponent under sample reuse is open. Ingredients exist: information-exponent results for online SGD \cite{benarous2021online}, early alignment \cite{boursier2022early}, and the mass clock and merging law of \cref{app:networks}.
\item \textbf{Composition costs, and width dependence at order 4.}
\item \textbf{Reservoir depletion.} When does the first escape use up the dormant units at practical scales?
\item \textbf{A bound of the form ``Occam for a few grown features plus a kernel norm for the background''}, as suggested by the hybrid model.
\item \textbf{Bayes versus gradient descent.} Characterize the targets on which their sample complexities are close.
\end{enumerate}
